\chapter{The survey: method}\label{chap:03-sondage-methode}

The survey, the first instrument of this report, rereads published papers with Pineau's reading grid,
extended by seven criteria drawn from her ICLR talk and by a second grid that we wrote. Both years
are read with the same grids by AI agents of the same model and every \enquote{yes} quotes a sentence
of the paper with its page. This chapter follows the order of the work: the protocol, the draws, the
grids, the readers and their double readings, the harmonisation and the checks, the verdict rule, the
statistics, then the re-readings added after the results and the deviations from the protocol.
Chapter~\ref{chap:04-sondage-resultats} gives the results.

\section{The protocol and its stages}\label{sec:sm-protocole}

\paragraph{A protocol written before the reading}
The survey protocol (\file{study/PROTOCOL.md}) was written before the first paper was
read. It fixes the question, the populations, the draws, the grid, the decision rules, the double
reading and the analysis. Every later change is entered at the bottom of the file, in a numbered list
of \enquote{deviations}. Each deviation is dated and states whether it was decided before or after the
result it concerns was known. The file holds 35 of them, numbered 1 to 36
(Section~\ref{sec:sm-ecarts}). The instruction sheet
given to the readers (\file{study/GRILLE.md}) summarises the protocol. The plan of the second grid
(\file{study/planification\_v2.md}), written before any coding, was made public as
it stood.

\paragraph{The stages}
Figures~\ref{fig:flux-sondage} and~\ref{fig:flux-sondage-2} show the complete flow, from the populations to the final counts. The
stages followed one another as below.
\begin{enumerate}
  \item \textbf{Protocol and grid v1.} Protocol, first draw of 50 papers from 2025 and 2026, calibration on 12
  papers from 2018, reading with grid v1 (Q1 to Q8, X1 to X7) by seven readers, one of them a second
  reader.
  \item \textbf{Grid v2 and enlargement.} Grid v2 (Y1 to Y12) and its decision rules, written before any coding
  (deviations 11 and 14). Enlargement of the two samples, written before any draw (deviation 12).
  Human rereading planned (deviation 13). Harmonisation of the 2018 calibration caught up (deviation
  15). Pilot (deviations 16 and 17), enlarged reading, audit of the raw readings, harmonisation
  (deviations 18 and 19), then first percentages. Same-grid verdict rule (deviation 20).
  \item \textbf{Robustness and trajectory.} Robustness analyses (deviation 21), checklist declarations (A5,
  deviation 22), human rereading abandoned (deviation 23), sensitivity to the title filter (A6,
  deviation 24), trajectory from 2019 to 2024 (A4, deviations 25 and 26).
  \item \textbf{Counter-checks.} Keyword search against the negative decisions (deviation 27), provenance
  of the readers (deviation 28), sensitivity to the change of model (deviation 29).
  \item \textbf{Technical audit and re-readings.} Response to a technical audit, each deviation written before its
  own computation: complementary analyses of the data already read (deviation 30), a re-reading of Q5
  on the same kind of version in both years (deviation 31), criterion X8 (deviation 32), a reader from
  another family of models (deviation 33). Then, after the results, a human
  rereading of both years, preregistered and not yet done (deviation 35), then eight fresh readers who
  each read both years (deviation 36).
\end{enumerate}

\fig{flux-sondage}{The survey, first part: from the populations to the samples. The three populations are the RL papers of ICML, NeurIPS and ICLR 2018, those of ICML and NeurIPS in 2019, 2020, 2022 and 2024, those of ICLR 2026, ICML 2026 and NeurIPS 2025. Every sample comes from a draw whose random seed is written in the protocol: 50 papers from 2018 (48 with experiments), 20 papers per year for 2019, 2020, 2022 and 2024, 50 classical RL papers from 2026. The first 2026 draw is kept as context.}

\fig{flux-sondage-2}{The survey, second part: from reading to comparisons. Every decision of the readers quotes a sentence of the paper with its page. Three sets of blind double readings, the harmonisation of the rules, the quote audit and the search for evidence against every \enquote{no} test these decisions before the final comparisons.}

\section{Populations and draws}\label{sec:sm-tirages}

\Cref{tab:echantillons} summarises, for each sample, the population, the draw and what enters the
analysis.

\paragraph{The definition of an RL paper}
A paper belongs to the population if it is accepted to the main track of a conference and if its
title matches, ignoring case, the regular expression
\begin{center}
\small\verb!reinforcement learning|\bRL\b|\bRLHF\b|\bRLVR\b|policy optimi[sz]ation|!\\
\small\verb!policy gradient|q-learning|actor-critic|\bMDPs?\b!
\end{center}
The same expression serves for every year. The three drawing scripts copy it verbatim:
\file{study/sample\_v2.py}, \file{study/trajectory\_sample.py} and \file{study/title\_filter.py}. A
filter on the title misses the RL papers whose title does not say so. Analysis A6 measures this loss
(Section~\ref{sec:sm-robustesse}).

\paragraph{The 2026 papers}
The three most recent conferences when the protocol was written are ICLR 2026, ICML 2026 and NeurIPS 2025.
Their public lists hold 17,836 papers, 1,252 of them on RL: 420 at ICLR, 494 at ICML, 338 at
NeurIPS. The list of RL papers is \file{study/rl\_population\_2026.json}. The first draw takes 17 papers
from ICLR, 17 from ICML and 16 from NeurIPS with \file{random.Random(2018)}, in that order, by
\file{sample()} on the list of each conference. The script \file{study/draw\_2026.py} redoes this draw
and checks that it recovers the same papers in the same order. We took the PDFs from the official
proceedings for ICLR and from arXiv for ICML and NeurIPS, since OpenReview refuses automated downloads. An
arXiv paper is accepted when its title is at least 85\% similar to the one in the list
(\file{difflib.SequenceMatcher}). Eight papers with no freely accessible PDF are replaced, in order, by
a \file{random.Random(2026)} draw from the rest of the same conference
(\file{study/replacements\_2026.json}). The final sample is \file{study/final\_2026.json}.

\paragraph{The initial calibration on 2018}
To check that our readers roughly recover Pineau's figures, the protocol adds 6 RL papers from ICML
2018 and 6 from NeurIPS 2018, with the same filter and the same random seed
(\file{study/calib\_2018.json}). The separate 2018 appendices are read: the PMLR supplementary material
for ICML, the proceedings archives for NeurIPS (deviation 8).

\paragraph{Enlargement A2: 50 papers from 2018}
With 12 papers from 2018 and 14 classical RL papers from 2026, the statistical audit of the first draw
(\file{study/audit\_stats.py}) lacks the power to detect a difference of 30 points. Deviation 12, written
before any draw and before any reading of the added papers, enlarges both samples. For
2018, the population brings together the three conferences Pineau names: 131 RL papers, 54 at ICML, 52
at NeurIPS and 25 at ICLR (\file{study/rl\_population\_2018.json}). The ICLR 2018 list, 336 papers of
the main track, orals and posters together, was taken from the official programme of the conference,
since the OpenReview and dblp interfaces refused automated requests at the time of the draw. The 50 papers
are allocated to conferences in proportion to their size, by the largest-remainder method: 21 for
ICML, 20 for NeurIPS, 9 for ICLR. The 12 initial papers stay. The 38 others (15, 14 and 9 per
conference) are taken in the order of a \file{random.Random(2018)} shuffle of the rest of the
population. A paper with no freely accessible PDF would be replaced by the next one in that order. No
replacement was needed (\file{study/replacements\_2018.json}). The PMLR and NeurIPS proceedings
supplied the PDFs, an Internet Archive copy of them when the proceedings site did not respond, or arXiv
by title. For ICLR 2018, OpenReview refusing automated downloads, the PDF is the Internet Archive
copy of the same OpenReview address, or else the arXiv version. The sample is
\file{study/calib\_2018\_v2.json}.

\paragraph{Enlargement A3: 50 classical RL papers from 2026}
The starting population is that of the 1,252 papers minus the 50 of the first draw, or 1,202
papers, shuffled with \file{random.Random(2026)}. A screen on the title alone, fixed in deviation 12,
skips the titles that contain \enquote{LLM}, \enquote{language model}, \enquote{reasoning},
\enquote{RLHF}, \enquote{RLVR}, \enquote{GRPO}, \enquote{preference}, \enquote{alignment},
\enquote{reward model}, \enquote{chain-of-thought}, \enquote{math} or \enquote{agentic}. The screen
sets aside 429 papers (\file{study/classic\_screened\_out\_2026.json}) and leaves an ordered queue of
773 papers (\file{study/classic\_queue\_2026.json}). The readers code the family of each paper in the
queue (criterion X5), in order, up to the thirty-sixth \file{classic} paper. The 9 ICLR 2026 papers at
the head of the queue were unreachable at the time of reading, the proceedings site being down: they are
skipped. The queue is frozen so that the selection does not depend on the reading date (deviation 17).
In all, 70 papers of the queue were read. The first 36 \file{classic} papers form the sample. Before
the thirty-sixth, 18 papers belonged to large language models (\file{llm}) and 3 to theory
(\file{theory}). The last 13 were read after the last paper retained (\cref{tab:articles-2026}). The
\file{llm} and \file{theory} papers stay coded and reported apart. They form no sample, since the
title screen set aside many others like them. The 2026 classical RL sample brings together the 14
\file{classic} papers of the first draw and these 36 papers. Deviation 18 refines X5 before this sample
is frozen: applied to the 120 papers read, the rule changes one code, K006, from \file{llm} to
\file{classic}.

\paragraph{Trajectory A4: 2019, 2020, 2022 and 2024}
Deviation 25, written before any draw, adds four years to date the change. The
population of each year brings together the RL papers of the main proceedings of ICML (PMLR volumes 97,
119, 162 and 235) and of NeurIPS. ICLR is missing, as OpenReview refuses automated downloads. After the
title screen of deviation 12, the population counts 134 papers in 2019, 222 in 2020, 301 in 2022 and
374 in 2024 (\file{study/trajectory\_queue.json}). The reading order is a
\file{random.Random(<year>)} shuffle of the two conferences pooled. The readers read the first 24 papers
of each year that have a PDF. The sample takes, per year, the first 20 \file{classic} papers with
experiments, in order. In 2019 and 2020, the first 24 papers held only 19 and 18 of them: the next
three papers of each year were read, as planned. Eight readers (M1 to M8) each read twelve papers taken
from all years, so that no reader is confounded with a year. M7 reads eighteen, including the six extra
papers of 2019 and 2020. A second pass was needed: the NeurIPS 2019 to 2022 appendices are zip
archives, which the first download had left aside. The raw readings of the first pass are recorded
in the repository (commit \file{0d3519a}). Each reader then rereads their papers with the appendix and
changes the decisions the appendix justifies.

\paragraph{Analyses A5 and A6}
Analysis A5 (deviation 22) takes the 30 NeurIPS 2025 papers already read, without a draw: 16 from the
first draw and 14 from the A3 queue. Analysis A6 (deviation 24) draws 400 papers, with
\file{random.Random(2029)}, among the papers of the three 2026 lists whose title does not match the
expression.

\begin{table}[tbp]
\centering\small
\caption{The samples of the survey. For each: the population, the drawing rule with its random seed,
what was read and what enters the analysis. Sources: the JSON files of \file{study/} named in the
text.}
\label{tab:echantillons}
\begin{tabular}{@{}>{\raggedright\arraybackslash}p{2.6cm} >{\raggedright\arraybackslash}p{4.1cm} >{\raggedright\arraybackslash}p{4.1cm} >{\raggedright\arraybackslash}p{3.2cm}@{}}
\toprule
Sample & Population & Draw & Analysed \\
\midrule
First draw 2026 & 1,252 RL papers of ICLR 2026, ICML 2026, NeurIPS 2025 & \file{Random(2018)}: 17, 17 and 16, then replacements by \file{Random(2026)} & 50 papers, 45 with experiments \\
Initial calibration 2018 & RL papers of ICML and NeurIPS 2018 & \file{Random(2018)}: 6 and 6 & 12 papers \\
A2, 2018 & 131 RL papers of ICML, NeurIPS and ICLR 2018 & quotas 21, 20 and 9, \file{Random(2018)} shuffle & 50 papers, 48 with experiments \\
A3, classical RL 2026 & 1,202 remaining papers, 773 after the title screen & \file{Random(2026)} shuffle, stop at 36 \file{classic} & 70 read, 14 + 36 = 50 retained \\
A4, 2019 to 2024 & 134, 222, 301 and 374 papers after the screen & \file{Random(<year>)} shuffle & 20 papers per year \\
A5, checklists & 30 NeurIPS 2025 papers already read & none & 30 papers \\
A6, title filter & 2026 papers outside the filter & \file{Random(2029)}: 400 & 400 abstracts, 20 papers read \\
\bottomrule
\end{tabular}
\end{table}

\section{The reading grids}\label{sec:sm-grilles}

\paragraph{Two versions of the grid}
Table~\ref{tab:criteres} lists the 27 criteria of the two grids, plus criterion X8, added after the
results (Section~\ref{sec:sm-x8}). Grid v1 brings together Pineau's eight criteria (Q1 to
Q8), as she presents them at NeurIPS (page 29), then seven criteria drawn from her ICLR talk (X1 to
X7). Grid v2 (Y1 to Y12) covers the flaws that v1 leaves unmeasured. It was written
(deviation 11), with its decision rules (deviation 14), before any v2 coding. Y11 and Y12 were added
next, after the search for related work, still before any coding.
Appendix~\ref{chap:A-grille} reproduces the rules in full.

\begin{small}
\begin{longtable}{@{}>{\raggedright\arraybackslash}p{0.9cm} >{\raggedright\arraybackslash}p{8.0cm} >{\raggedright\arraybackslash}p{2.2cm} >{\raggedright\arraybackslash}p{2.4cm}@{}}
\caption{The reading criteria. The \enquote{Values} column gives the permitted answers. A
\file{yes}/\file{no}/\file{na} criterion takes \file{na} when it does not apply. The last column gives
the flaw of Table~\ref{tab:failles} that the criterion measures, or its role.}\label{tab:criteres}\\
\toprule
Code & Criterion & Values & Flaw or role \\
\midrule
\endfirsthead
\toprule
Code & Criterion & Values & Flaw or role \\
\midrule
\endhead
\bottomrule
\endfoot
\multicolumn{4}{@{}l}{\emph{Grid v1: Pineau's criteria}} \\
Q1 & The paper has experiments & yes, no & scope \\
Q2 & It uses neural networks & yes, no, na & population \\
Q3 & All of the method's hyperparameters are given & yes, no, na & F2, F4 \\
Q4 & All of the baselines' hyperparameters are given & yes, no, na & F7 \\
Q5 & The code is linked & yes, no, na & F1 \\
Q6 & How the hyperparameters were chosen is described & yes, no, na & F5 \\
Q7 & The evaluation uses held-out conditions & yes, no, na & F15 \\
Q8 & A statistical test is applied to a comparison & yes, no, na & F12 \\
\midrule
\multicolumn{4}{@{}l}{\emph{Grid v1: our additions drawn from the ICLR talk}} \\
X1 & Number of independent training runs behind the main results & integer & F9 \\
X2 & An uncertainty is reported on the main results & yes, no, na & F12 \\
X3 & The best runs are selected & yes, no, na & F10 \\
X4 & The paper says it tuned the baselines with the same care & yes, no, na & F7 \\
X5 & Family: \file{llm}, \file{classic} or \file{theory} & one of the three & stratum \\
X6 & For \file{llm}: number of base-model families & integer & F16 \\
X7 & The compute used is quantified & yes, no, na & F2, F3 \\
\midrule
\multicolumn{4}{@{}l}{\emph{Grid v2}} \\
Y1 & The performance measure is defined & yes, no, na & F8 \\
Y2 & The error bars are defined & yes, no, na & F8 \\
Y3 & Number of tasks behind the main conclusion & integer & F13 \\
Y4 & A task where the method does not win is shown & yes, no, na & F13 \\
Y5 & The implementation of the baselines is identified & yes, no, na & F14 \\
Y6 & Exclusions are stated, or their absence & yes, no, na & checklist \\
Y7 & The version of the environments or benchmarks is given & yes, no, na & F16 \\
Y8 & A generalisation test outside the training conditions & yes, no, na & F15 \\
Y9 & The sensitivity to at least one hyperparameter is shown & yes, no, na & F6 \\
Y10 & The software dependencies are listed with their versions & yes, no, na & F3 \\
Y11 & The type of uncertainty is named & yes, no, na & F12 \\
Y12 & A statistical tool later than 2018 is used & yes, no, na & F12 \\
\midrule
\multicolumn{4}{@{}l}{\emph{Added after the results (deviation 32)}} \\
X8 & The reported uncertainty is computed over independent random seeds & yes, no & F11, F12 \\
\end{longtable}
\end{small}

\paragraph{Decision rules of grid v1}
The protocol fixes the following rules, completed by deviations 1 to 6.
\begin{itemize}
  \item \textbf{Q3 and Q4.} \file{yes} if a table or a paragraph gives the hyperparameters of all the
  main methods. \enquote{Default values of [library]} without listing them counts as \file{yes} only
  if the library version is given. Q4 is \file{na} when no baseline is compared.
  \item \textbf{Q5.} A link to code in the paper. \enquote{On request}, \enquote{will be released} and
  \enquote{will be made public at <URL>} count as \file{no}. A project page counts as \file{yes} only
  if it leads to code, which was checked by visiting each page (deviations 6 and 19).
  \item \textbf{Q6.} The selection method is stated: grid, random search, budget, criterion.
  \enquote{Tuned} and nothing more counts as \file{no}.
  \item \textbf{Q7.} Evaluation on tasks, environments, environment random seeds or datasets unseen in
  training and in tuning. For \file{llm}, benchmarks distinct from the RL training data count as
  \file{yes}. A \file{theory} paper with toy experiments counts as \file{no} (deviation 4).
  \item \textbf{Q8.} A test applied to a comparison: $t$, Welch, Mann-Whitney, bootstrap with an
  interval, rliable IQM with stratified intervals. Error bars with no test count as \file{no} for Q8
  and as \file{yes} for X2. X2 therefore counts any error bar, whatever its source. Criterion X8,
  added after the results, asks which source (Section~\ref{sec:sm-x8}).
  \item \textbf{X1.} The smallest number of independent training runs stated for a main table or
  curve. For \file{llm}, evaluation samples (pass@$k$, \enquote{16 samples}) do not count.
  \item \textbf{X3.} \file{yes} also when the paper's tables show that the main result is the best of
  several identical training runs, even if the text does not say so (deviation 5).
  \item \textbf{X4.} \file{yes} requires a described tuning of the baselines or an announced equal
  tuning budget. \enquote{Same configuration as the method}, with no tuning described for anyone,
  counts as \file{no} (deviation 1).
  \item \textbf{X5.} The family is read from the model that the RL trains (deviation 18): \file{llm}
  if the optimised policy is a pretrained generative model conditioned on language, \file{classic}
  otherwise, even in the service of a large model. \file{theory} is reserved for papers with no
  experiment or with toy experiments only.
  \item \textbf{X7.} \file{yes} requires at least one quantity: hours, GPU-hours, FLOP, or a number of
  GPUs with a duration. The name of the hardware, with no quantity, counts as \file{no} (deviation 2).
\end{itemize}

\paragraph{Decision rules of grid v2}
Deviation 14 fixes the rules, deviation 17 refines them after the pilot, deviation 19 after the audit
of the raw readings.
\begin{itemize}
  \item \textbf{Scope.} If Q1 is \file{no}, all the other criteria are \file{na}. Y2 and Y11 are
  \file{na} if and only if X2 is \file{no}. Y4 is \file{na} if Y3 is 1. Y5 is \file{na} if and only if
  Q4 is \file{na}.
  \item \textbf{Y1.} The paper states what the reported number contains and at what point it is
  taken: final return, mean over the last $N$ episodes, best point, area under the curve, pass@$k$
  with $k$ and the number of samples. \enquote{Return} or \enquote{accuracy} and nothing more counts as
  \file{no}. A curve whose evaluation points are stated counts as a measurement point. For a quantity
  other than the score of a policy, the defined quantity suffices (deviation 19).
  \item \textbf{Y2.} The paper states what the band or the interval represents. Bars present but
  undefined count as \file{no}.
  \item \textbf{Y3.} An integer, the number of distinct training settings behind the main conclusion.
  Two difficulty levels of the same game, or the quality variants of an offline dataset, count as one.
  A meta-training counts as one. A setting present only in an appendix does not count. The analysis
  turns Y3 into \file{yes} when the integer is at least 2.
  \item \textbf{Y4.} A task-by-task result shows at least one task where the method does not win.
  \item \textbf{Y5.} For each main baseline, the origin of the code is stated: repository, named
  library, \enquote{reimplemented by us}, numbers copied with their source, released model designated
  by its exact name. An implied reimplementation does not count.
  \item \textbf{Y6.} The runs, random seeds or data removed are stated, or their absence is stated.
  Silence counts as \file{no}.
  \item \textbf{Y7.} The version of each main environment or benchmark: Gymnasium suffix, suite
  version, name that fixes a single published set. \enquote{MuJoCo} without a version counts as
  \file{no}. \enquote{MATH-500} counts as \file{yes}, bare \enquote{MATH} as \file{no}.
  \item \textbf{Y8.} At least one evaluation on conditions absent from training and tuning. Stricter
  than Q7: for \file{llm}, a different benchmark counts only if it is presented as an
  out-of-distribution test. New training random seeds are not held-out conditions.
  \item \textbf{Y9.} At least one hyperparameter of the method is swept with the result shown. A
  component ablation does not count.
  \item \textbf{Y10.} The versions of the main libraries are given, or an environment file is cited.
  The repository is not visited, unless the paper points explicitly to the file.
  \item \textbf{Y11.} The type of uncertainty is named somewhere, with a \file{kind} field
  (\file{std}, \file{se}, \file{ci}, \file{iqr}, \file{minmax}, \file{other}).
  \item \textbf{Y12.} At least one tool among the IQM, the stratified bootstrap, rliable, performance
  profiles, the probability of improvement. A plain bootstrap that yields a reported interval counts.
\end{itemize}

\paragraph{What a share measures}
For a criterion and a sample, the share is the number of \file{yes} divided by the number of
\file{yes} and \file{no}. Papers with no experiment (Q1 \file{no}) and \file{na} values leave the
denominator (\file{tally()} in \file{src/pineau\_rl/grid.py}). A denominator smaller than the sample
size therefore signals papers where the criterion does not apply, for instance a paper with no compared
baseline for Q4. Pineau reports 100\% experimental papers: her denominators are the 50 papers.

\section{The readers}\label{sec:sm-lecteurs}

\paragraph{Agentic AI models}
The reading is entrusted to agentic AI models, generative AI models that act: they open the PDF, find the sentence behind each decision and fill in the grid. We also call them AI agents. They are Claude Code subagents, each launched with a written
instruction, a list of identifiers and a folder to write in. The folder \file{study/readers/} gives the
provenance of each agent, reconstructed from the session transcripts (deviation
28): the exact instruction (\file{prompts/<reader>.md}), the model recorded in the transcript, the
dates, the papers written and the way the PDFs were read (\file{readers.csv}). The 38 reading folders
each match one agent and its papers, without discrepancy (\file{folder\_check.csv}). All readers ran
on the model \file{claude-opus-5-5}, except two that finished on \file{claude-opus-5} after a content
filter of the provider interrupted an answer. The two are C2 (papers ML14 to ML17) and M4 (nine A4
papers), thirteen papers in all.

\paragraph{The work of a reader}
Each reader reads the whole paper, appendices included. The reader converts the PDF to text with
\file{pdftotext} and often looks at the page images. The reader writes one JSON file per paper, in
\file{study/codes/<reader>/} for grid v1 and \file{study/codes\_v2/<reader>/} for v2. For each
criterion, the file gives the value, the page and an exact quotation of the paper copied word for
word. If nothing justifies a \file{yes}, the quotation is \enquote{absent after reading the
experiments, appendices and links sections} and the page is \file{null}. The reader records hesitations
in a \file{doubts} field. A reader never reads another reader's codes, nor the results, nor the
documents of \file{docs/}. For a paper already read in v1, the v2 reader receives only the v1 values of
Q1, Q4 and X2, which fix the scope rules.

\paragraph{The mechanical check}
The script \file{study/check\_codes.py} checks each v2 file before any percentage. It checks the schema
(criteria present, permitted values, integers, \file{kind} field of Y11) and the scope rules
(\file{validate()} in \file{src/pineau\_rl/grid.py}). It then searches for each quotation in the text of
the PDF and of its appendix. The reader's notes in square brackets are removed. A quotation stitched from
pieces joined by \enquote{[...]} is cut into its pieces. A piece of fewer than 12 letters is too
short to test. The comparison uses letters and digits only, so that spaces, hyphenation, ligatures and
mathematical symbols do not count. A piece not found verbatim is searched for approximately: it passes
if 90\% of its characters are found in a window of the text. A failing file goes back to its reader. A
second failure on the same criterion is decided from the PDF and recorded in
\file{study/overrides\_v2.json}.

\paragraph{Three facts outside the protocol}
The reconstruction of deviation 28 established three facts the protocol left unstated. Readers C2 and
M4 finished on the model \file{claude-opus-5}. The primary reader M7 opened the file K002 of reader K1
as a format example. K002 is not M7's and enters no double reading. The grid v1 readings (readers A to
G) took place in a session opened in a folder other than the repository. The reconstruction
also checks that the second readers G, G2, R2 and R3 opened no file of another reader.

\paragraph{In the main comparison, year and readers are confounded}\label{sec:sm-annees}
No primary reader read both years of the main comparison. The papers of 2018 come from readers F, F2
and H1 to H4, those of classical RL in 2026 from readers A to E, A2 to E2 and K1 to K6
(\file{study/survey\_extra.json}, key \file{readers\_in\_both\_years}, empty). All are copies of the
same model. Each one nonetheless has its own session, its own written instruction and its own date. An
effect proper to the readers of one year therefore mixes with the effect of the year. The cut that
removes one reader at a time (Section~\ref{sec:sm-robustesse}) does not see it. Deviation 30 (d),
written before its computation, measures what the data allow us to see.
\begin{itemize}
  \item \textbf{Within a year}, the share of a criterion varies from one reader to another. A
  permutation test measures it: the chi-square statistic of the readers $\times$ answers table,
  \num{10000} permutations of the answers between readers, NumPy generator with seed 2031, readers of
  at least five papers. In 2018, Q5 ranges from 0 out of 8 (H3) to 6 out of 11 (F), $p = 0.049$. In
  2026, X7 ranges from 17\% to 83\%, $p = 0.027$. The six other (criterion, year) pairs behind the
  verdicts have $p > 0.18$. These differences between readers mix the reader with the batch of papers: F read the twelve
  papers of the initial calibration, eleven of them with experiments (the twelfth, C23, has none and
  leaves the count), the H readers the thirty-eight added ones.
  \item \textbf{Between the years}, the only reader who read both is R2, on five papers of each. On
  these ten papers, R2's difference between 2026 and 2018 equals the one the primary readers obtain on
  the same ten papers, for the four criteria behind the verdicts. For X2, this difference is a fall
  ($-0.2$: 5 of 5 in 2018, 4 of 5 in 2026), whereas the X2 verdict, over the 98 papers, is a rise. Five papers per year do not close the question.
\end{itemize}
The readers of criterion X8 already each read both years (Section~\ref{sec:sm-x8}). Deviation 36 does
the same for the four criteria behind the verdicts: eight fresh readers reread the 98 experimental
papers, each of them on both years (Section~\ref{sec:sm-doubles}).

\section{The blind double readings}\label{sec:sm-doubles}

\paragraph{Four second readers}
A second reader rereads papers without seeing the first reading or knowing who made it.
\begin{itemize}
  \item \textbf{G (grid v1) and G2 (grid v2).} The first 10 papers of the first draw in identifier
  order after a \file{random.Random(7)} shuffle.
  \item \textbf{R2 (both grids).} 10 papers drawn by \file{random.Random(7)} among the papers added by
  A2 and A3: 5 from 2018 and 5 from 2026.
  \item \textbf{R3 (both grids).} 10 papers drawn by \file{random.Random(2030)} among the 80 of A4,
  reread with the appendices (deviation 26).
\end{itemize}
Thirty papers are reread in all. Only R2's five papers from 2018 measure agreement on papers from 2018,
before the complete rereading of deviation 36.
Since all readers are copies of the same model, agreement measures the stability of that model more
than the agreement of two independent readers (deviation 28). Two complete rereadings, written after
the results, answer this: eight fresh readers of the same model who each read both years (deviation
36), then a reader from another family of models (deviation 33, Section~\ref{sec:sm-gemma}).

\paragraph{Eight fresh readers on both years (deviation 36)}
Deviation 36, written after the results and before any reading, adds a complete
rereading of the 98 experimental papers of the main comparison, 48 from 2018 and 50 from 2026, on the
four criteria behind the verdicts: Q5, Y9, X7 and X2.
\begin{itemize}
  \item \textbf{The readers.} Eight fresh AI agents, T1 to T8, from the same model family as the
  original readers, had read no paper of the survey. Each receives its list, the path to the PDFs and the prompt
  \file{study/r36\_prompt.txt}, which copies out the rules of the four criteria. It sees neither the other readings nor the results. The year shows in every PDF: the readers
  get the rules, not the hypothesis.
  \item \textbf{The allocation.} A \file{random.Random(2035)} generator deals the papers of each year
  in turn, then shuffles each list so that the years alternate in the reading order
  (\file{study/r36\_assign.py}). The eight each read six papers from 2018 and six or seven from 2026.
  Year and reader are thus crossed.
  \item \textbf{The check.} Each quotation is looked up in the PDF as for X8
  (\file{study/r36\_check.py}). All 98 files pass, at once or after their reader's correction: no
  decision is invalid. A Q5 \enquote{yes} that points to a project page takes the final value of the
  original reading when that reading cites the same address, which holds for all eleven such \enquote{yes}. A variant of the analysis, the least favourable to the Q5 verdict, counts every project page as \enquote{yes}.
  \item \textbf{The analysis.} The script \file{study/r36\_stats.py} redoes the main comparison with
  these readings only: Newcombe, Fisher, Holm over the four criteria. It then fits, per criterion, a
  logistic mixed model. The probability of a \enquote{yes} depends on the year and on a random reader
  effect, normal with standard deviation $s$, fitted by maximum likelihood with 40-node Gauss-Hermite
  quadrature. A likelihood-ratio test judges the year effect, with Holm over the four criteria.
  Beside it, the agreement of the new readings with the final codes is computed per criterion and per
  year, with Fisher's test of the 2018 agreement against the 2026 one.
  \item \textbf{The rule, fixed in advance.} A verdict holds if the difference between years and the
  year effect in the model are both positive and if their two tests, Fisher's and the likelihood-ratio
  test, both pass the Holm correction at 5\%. It is \enquote{same
  direction, not established} if both are positive but one of the two tests fails, \enquote{not
  reproduced} otherwise.
\end{itemize}
Every paper already has its original reading. The new readings therefore form a complete double
reading with them, with no third reading. The results are in Section~\ref{sec:sr-annees}.

\paragraph{A reader from another family of models (deviation 33)}\label{sec:sm-gemma}
All our readers, the eight fresh ones included, belong to the same family of models. Deviation 33,
written after the results and before any reading, adds a reader from another family,
\texttt{gemma3:12b} (Google), served locally by Ollama at temperature 0, seed 0 and a context of
\num{98304} tokens. It asks two questions. Does this reader find the differences between years on Q5,
Y9, X7 and X2? Does its disagreement with our readers depend on the year?
\begin{itemize}
  \item \textbf{The sample and the input.} The 98 papers with experiments of the main comparison, in
  an order shuffled by \file{random.Random(2034)}. The text comes from \file{pdftotext}, reference
  section removed, truncated at \num{320000} characters. The prompt copies the rules of
  \file{GRILLE-en.md} for the four criteria, with their clarifications (deviations 2, 6, 14, 17 and
  19). One call per criterion (\file{study/gemma\_reader.py}).
  \item \textbf{The check.} A \enquote{yes} is valid only if its quotation is found in the PDF and
  contains its criterion's pattern, fixed for the search against the \enquote{no} decisions (deviation
  27). Otherwise there is a single retry, after which the decision is \enquote{invalid} and leaves the
  denominators. A malformed answer counts as invalid for its own criterion only. A \enquote{no} is
  always valid. The call per criterion and the pattern rule were fixed after a pilot of three papers,
  outside the analysis, before the full reading; the rule on malformed answers, before the
  rereading of the one paper concerned.
  \item \textbf{The analysis.} The script \file{study/gemma\_stats.py} gives, per criterion and per
  year, the agreement of the valid decisions with our final codes, with Wilson and a bootstrap kappa,
  then a Fisher test of the 2018 agreement against the 2026 agreement. It then recomputes the 2026
  minus 2018 difference from these decisions alone, Newcombe and Fisher with Holm over the four
  criteria, next to our readers' difference on the same papers. An announced variant counts every
  invalid decision as \enquote{no}.
  \item \textbf{The rule, fixed in advance.} Same sign and overlapping intervals:
  \enquote{reproduced}. Same sign and this reader's interval containing zero: \enquote{same
  direction, not established}. Opposite sign: \enquote{not reproduced}. The protocol notes that a
  model of 12 billion parameters reads less well than a large model: a disagreement does not by itself
  refute our readings, while a disagreement that depends on the year would be the bias sought.
\end{itemize}
The results are in Section~\ref{sec:sr-gemma}.

\paragraph{Agreement, on the raw readings}
Agreement is computed on the raw readings, before any harmonisation: harmonisation applies to both
readings and would inflate agreement. For each criterion, \file{study/agree\_v2.py} keeps the papers
where both readers answered \file{yes} or \file{no} (Y3 transformed as in the analysis), then computes
the share of identical answers and Cohen's kappa when both readers vary. A criterion below 70\%
agreement is called fragile. The agreement of grid v1 on the first draw is computed by
\file{study/audit\_stats.py}, which also counts the \file{na} values.

\paragraph{Intervals on agreement}
On ten papers, an agreement of 9 out of 10 has the Wilson interval $[60,\,98]$\%: one pair of
readers says little. Deviation 30 (a), written before its computation, adds two intervals
(\file{study/survey\_extra.py}). The share of agreement receives its Wilson interval
(Section~\ref{sec:sm-stats}). Kappa receives a percentile bootstrap interval over the papers:
\num{10000} resamples with replacement, NumPy generator with seed 2030. A resample where kappa is
undefined is dropped and counted. The three pairs are then pooled by criterion: G for grid v1 or G2
for v2, with R2 and R3, that is at most thirty papers per criterion
(results in Section~\ref{sec:sr-accord}).

\paragraph{The human rereading, abandoned then preregistered again}
Deviation 13 planned for us to reread 10 papers from 2026 ourselves, alone, on seven criteria,
without seeing the agents' readings. The papers had been drawn by \file{random.Random(2027)}. Deviation
23 abandons this rereading before it started. The folder
\file{study/relecture\_humaine/} remains as a trace of the plan. No human reread the papers
blind. Each decision remains checkable through its quotation and its page.

Deviation 35, written after the results and before any draw, plans a new human
rereading, this time on both years. We will reread thirty papers of the main comparison ourselves, alone,
fifteen per year, drawn by \file{random.Random(2034)} among those that no double reading and no part of
the deviation 13 plan had already covered.
\begin{itemize}
  \item \textbf{Blindness.} Each paper gets a sheet with an opaque name (H01 to H30) and a renamed copy of its PDF (\file{study/relecture\_humaine/feuilles/}). The sheet shows no identifier, title,
  conference, year or agent value. The PDF still shows the year. We also know the hypothesis: the
  deviation says so in advance.
  \item \textbf{Criteria.} Q5, Y9, X7, X2 and X8, with the agents' rules copied into each sheet. Every
  decision comes with an exact quotation and its page.
  \item \textbf{Analysis.} The script \file{compare.py} refuses to run while any of the thirty sheets is
  incomplete. It will then read the agents' codes, measure the human's agreement with them and say
  whether the difference between years holds under a human reading.
\end{itemize}
The thirty sheets are still empty.

\section{The pilot and the harmonisation}\label{sec:sm-harmonisation}

\paragraph{The pilot}
Before the primary readers, a pilot reader codes three papers (deviation 16): LR08 (\file{llm}) and
ML15 (\file{classic}), chosen by \file{random.Random(11)} among the 40 papers from 2026 outside the
double reading, then C218, the first of the 38 new papers from 2018 after a \file{random.Random(11)}
shuffle. The pilot serves only to find ambiguous rules and to test \file{check\_codes.py}. Its files
(\file{codes\_v2/pilot/}) enter no analysis. The three papers are recoded by their regular reader.
Deviation 17 writes down the resulting clarifications, before the main coding and before any v2
percentage.

\paragraph{The harmonisation of grid v1}
Deviations 1 to 6 were written during the coding of the first draw, before aggregation. They are
applied from the quotations, in \file{study/overrides.json}. The check of the quotations not found
showed that these rules had been applied only to the 2026 papers (deviation 15). The
comparison between years therefore mixed two slightly different grids. The same rules were applied to
the 12 papers from 2018, with no new rule. For X4, C11, C13, C14 and C16 switch to \file{no}. For Q7,
C14 switches from \file{na} to \file{no}. For X7, C14 switches to \file{no}. The script
\file{study/analyze\_final.py} now regenerates \file{study/results\_final.json}, until then written
without a script. This correction was decided after seeing the percentages, which the protocol
declares.

\paragraph{The harmonisation of grid v2}
After the enlarged reading, an audit rereads all the v2 files of the primary readers, criterion by
criterion, without computing any share (\file{study/harmonisation\_audit.md}, 170 files). It finds
codes contrary to the rules and ten cases that no rule covered. Deviation 19, written
before any percentage, rules on these cases for both years. Its decisions are applied from the
quotations, each with its reason, in \file{study/overrides\_v2.json}: 80 changes, of which Y1 21, Y3 5
(and Y4 5 through the scope rule), Y5 10, Y6 3, Y7 19, Y8 6, Y9 9 and Q5 2. The project pages given as
a link to the code were visited. Five lead to the code, two do not. The link of K090
leads to an expired anonymous mirror (HTTP 410). The link of C19 leads to a library released after the
paper, without its scripts. We, the authors of the protocol, ruled on two cases, NE16 and ML14, all of whose
experiments are in the appendix, so that the rule on settings present only in an appendix does not
apply to them.

\paragraph{Three criteria with no new rule}
For X3, X4 and X7, the audit found borderline cases, such as compute reduced to inference time or
failed random seeds set aside. We added no new rule. A rule written at that point would have applied
only to the papers of the enlarged reading, which deviation 15 forbids: the two years would again have
been read with two different grids. These cases stay coded as the reader read them.

\paragraph{The final value of a criterion}
The function \file{value()} of \file{src/pineau\_rl/grid.py} computes the value that enters the
analysis. It takes the v2 reading of the primary reader when that reading contains the criterion, with
\file{overrides\_v2.json}. Otherwise it takes the v1 reading, with \file{overrides.json}. The readings of
the second readers and of the pilot never enter the analysis.

\section{The checks of quotations and of negative decisions}\label{sec:sm-controles}

\paragraph{The audit of the grid v1 quotations}
The first pass of the automatic check on the v1 files left 89 quotations not found. An audit agent
checked all of them in the PDFs (\file{study/v1\_quotes\_audit-en.md}). None is invented: 83 are present
word for word, 1 is a faithful paraphrase, 5 have a value contrary to the rules. The failures had three causes:
French comments that the reader had added to an exact English quotation, the two-column layout that
\file{pdftotext} interleaves, three NeurIPS 2018 appendices missing from the \file{study/pdf/}
folder. The five values contrary to the rules led to deviation 15.

\paragraph{The keyword search against the \enquote{no} decisions}
A \file{yes} decision rests on a quotation and a page. A \file{no} decision rests on no quotation. Yet the
\enquote{changed} verdicts rest mostly on the \file{no} decisions of 2018. Deviation 27, written
after the results and before any search, checks every final \file{no} decision of the
two samples of the main comparison against the text of the paper, on the seven criteria Q4, Q5, Q7,
Q8, X2, X7 and Y9. The script \file{study/falsify\_no.py} extracts the text of the paper and of its appendix, removes
the references section, then searches for patterns fixed in the deviation and never modified
afterwards. For example, the pattern of Q5 is
\begin{center}
\small\verb!github|gitlab|bitbucket|anonymous\.4open|code (is |are |will be )?!\\
\small\verb!(publicly )?available|source code|our (code|implementation)|open[- ]source!
\end{center}
Each match is kept with 300 characters of context. An arbiter, an agent that read neither sample and
sees neither the paper nor the year, says according to the rules of \file{GRILLE.md} whether the
excerpt suffices to turn the decision into \file{yes}, or whether it remains doubtful. The batches the
arbiter saw and its raw answers are in \file{study/falsify\_batches/}. The script
\file{study/falsify\_stats.py} then redoes the main comparison with the turned decisions, without
changing the test, as the deviation announces. It also redoes the comparison with the turned decisions
plus the doubtful ones, a bound the deviation did not announce.

\section{The calibration on 2018 and the verdict rule}\label{sec:sm-verdict}

\paragraph{The reason for rereading 2018}
Pineau describes neither her draw nor how she decided doubtful cases. Comparing our 2026 readings with
her figures would mix two changes: the change in practice and the change of grid. The calibration
measures the second. From the first 12 papers on, our grid looked stricter than hers on Q3 and more
lenient on Q7 (deviation 7). The protocol then planned to conclude that a practice had changed only if it held both
against her figures and against our own reading of 2018.

\paragraph{The same-grid verdict rule}
With 50 papers from 2018, the secondary comparison of deviation 12 confirmed that our grid reads 2018
differently from Pineau (Chapter~\ref{chap:04-sondage-resultats}). Deviation 20, decided
after these first results, draws from it the rule of the report: a verdict rests only on the main
comparison of deviation 12, our 50 classical RL papers from 2026 against our reading of 50 papers from
2018, with the same grids, read by agents of the same model. This comparison had been fixed before
any reading of the added papers. The change of rule, by contrast, followed the results, which the
protocol declares. The comparisons with Pineau's figures stay on display as context. A criterion
\enquote{changed} when the Fisher $p$-value of its difference between years, adjusted by the Holm
correction, stays below 5\%. We then say that the difference passes the correction. The agents of
2018 are not those of 2026: in this comparison the year effect is confounded with the readers'
(Section~\ref{sec:sm-annees}), which deviation 36 separates (Section~\ref{sec:sr-annees}).

\paragraph{What the original rule would have given}
Deviation 30 (c) applies the rule of deviation 7 to the final data. It was written before the
repository's computation, but after the technical review had sketched its result
(Section~\ref{sec:ecarts-audit}). This rule applies to Q3 to Q8, the criteria for which Pineau gives a
figure. A criterion \enquote{changed} under it when the difference between classical RL in 2026 and
Pineau's figures passes the Holm correction over these six criteria. The change is \enquote{confirmed}
when the difference from our reading of 2018 passes it too, in the same direction. It is
\enquote{ambiguous} when the difference passes only against our reading of 2018. X2, X7 and Y9 have no
figure from Pineau: the original rule did not judge them. The result, in Section~\ref{sec:sr-pineau}
with \cref{tab:regle-origine}, makes the change of rule checkable, a change the calibration on Q3
justifies (55\% in our reading of 2018, 90\% in Pineau's).

\section{The statistics}\label{sec:sm-stats}

All the statistical functions are in \file{src/pineau\_rl/stats.py}. They call SciPy and statsmodels
and are deterministic. Six scripts use them: \file{study/extension\_stats.py} (main comparison
and secondary comparisons), \file{study/trajectory.py} (A4), \file{study/agree\_v2.py} and
\file{study/audit\_stats.py} (agreement), \file{study/robustness.py} (robustness) and
\file{study/falsify\_stats.py} (negative decisions). In what follows, $k$ is the number of \file{yes},
$n$ the denominator and $\hat{p} = k / n$.

\paragraph{Wilson interval}
The 95\% interval of a share is that of \textcite{wilson1927probable}, suited to small samples and to
shares close to 0 or 1. With $z = \Phi^{-1}(0.975)$,
\begin{equation}
  \frac{\hat{p} + \frac{z^2}{2n}}{1 + \frac{z^2}{n}}
  \;\pm\;
  \frac{z}{1 + \frac{z^2}{n}} \sqrt{\frac{\hat{p}(1 - \hat{p})}{n} + \frac{z^2}{4n^2}} .
\end{equation}

\paragraph{Newcombe interval for a difference}
Pineau's 50 papers also form a sample. Treating her shares as exact, the 55\% of papers with code for
instance, would overstate the certainty of the verdicts. The difference $d = \hat{p}_a - \hat{p}_b$ between
two independent samples therefore receives the hybrid interval of \textcite{newcombe1998interval} (his
method 10). With the Wilson intervals $(l_a, u_a)$ and $(l_b, u_b)$ of each share,
\begin{equation}
  \Bigl[\, d - \sqrt{(\hat{p}_a - l_a)^2 + (u_b - \hat{p}_b)^2} ,\;
  d + \sqrt{(u_a - \hat{p}_a)^2 + (\hat{p}_b - l_b)^2} \,\Bigr].
\end{equation}
The computation is that of \file{confint\_proportions\_2indep(method="newcomb")} in statsmodels. For one
of Pineau's figures, the share enters the computation as $k = \mathrm{round}(50\,p)$ out of $n = 50$,
because a count is whole. For code (Q5), Pineau's 55\% thus becomes 28 papers out of 50, that is 56\%.
The text keeps her published figure, 55\%.

\paragraph{Fisher's exact test}
The $p$-value of a comparison is that of the exact test of \textcite{fisher1922interpretation} on the
table $\bigl(\begin{smallmatrix} k_a & n_a - k_a \\ k_b & n_b - k_b \end{smallmatrix}\bigr)$. Under the
null hypothesis and with fixed margins, $k_a$ follows a hypergeometric distribution. The two-sided
$p$-value sums the probabilities of the tables as probable as the observed one or less probable
(\file{scipy.stats.fisher\_exact}).

\paragraph{Holm correction}
Testing several criteria at once multiplies the chances of crossing the threshold by chance. The
correction of \textcite{holm1979simple} controls the probability of at least one false positive in a
family of $m$ tests. With the sorted $p$-values $p_{(1)} \leq \dots \leq p_{(m)}$, the adjusted value is
\begin{equation}
  \tilde{p}_{(i)} = \min\Bigl(1,\; \max_{j \leq i}\, (m - j + 1)\, p_{(j)} \Bigr).
\end{equation}
The families are fixed in the scripts. The main comparison has two: the 11 tested v1 criteria (Q2 to
Q8, X2, X3, X4, X7) and, separately, the 12 v2 criteria. The secondary comparisons, 2026 against Pineau
and our 2018 against Pineau, each take as their family the six criteria Q3 to Q8 for which Pineau gives
a comparable figure. For the trajectory, each family brings together the five tested criteria (Q5, Y9,
X7, X2, Q8).

\paragraph{Minimum detectable difference}
A difference that does not pass the threshold does not prove the absence of change. To say what the survey
could see, deviation 30 (b), written before the repository's computation (the technical review had
sketched its order of magnitude, Section~\ref{sec:ecarts-audit}), computes the exact power of the
two-sided Fisher test. Let $p_0$ be the share observed in our reading of 2018, $n_0$ and $n_1$ the
actual denominators of 2018 and 2026, $\alpha$ the threshold. If the true share in 2026 is $p_1$, the
power is
\begin{equation}
  \pi(p_1) = \sum_{k_1 = 0}^{n_1} \sum_{k_0 = 0}^{n_0}
  \binom{n_1}{k_1} p_1^{k_1} (1 - p_1)^{n_1 - k_1}
  \binom{n_0}{k_0} p_0^{k_0} (1 - p_0)^{n_0 - k_0}\;
  \mathbb{1}\bigl[\, p_{\text{F}}(k_1, k_0) < \alpha \,\bigr],
\end{equation}
where $p_{\text{F}}(k_1, k_0)$ is the Fisher $p$-value of the observed table. The minimum detectable
rise is the smallest $\delta$, in steps of one point, such that $\pi(p_0 + \delta) \geq 0.8$. The
minimum detectable fall is defined in the same way with $p_0 - \delta$. The computation runs at
$\alpha = 0.05$, then at the Bonferroni threshold $0.05/m$, with $m = 11$ for the v1 criteria and
$m = 12$ for the v2 criteria (\file{study/survey\_extra.py}). The same formula, with $p_1$ equal to the
share observed in 2026, gives the power the test had for the difference actually seen
(Section~\ref{sec:sr-principale}).

\paragraph{Cochran-Armitage trend test}
For trajectory A4, the test of \textcite{cochran1954some} and \textcite{armitage1955tests} looks for a
linear trend of the share over the years. With $K = \sum_i k_i$, $N = \sum_i n_i$, $\bar{p} = K / N$,
the scores $x_i$ equal to the years (2018, 2019, 2020, 2022, 2024, 2026) and
$\bar{x} = \sum_i x_i n_i / N$,
\begin{equation}
  T = \sum_i x_i\,(k_i - n_i \bar{p}),
  \qquad
  \operatorname{Var}(T) = \bar{p}(1 - \bar{p}) \sum_i n_i (x_i - \bar{x})^2,
  \qquad
  z = T / \sqrt{\operatorname{Var}(T)}.
\end{equation}
The two-sided $p$-value is read from the normal distribution. Two differences complete this test, each
with its Newcombe interval and a Fisher test: 2019 and 2020 pooled against 2018 (the immediate effect of the checklist), then 2022
and 2024 pooled against 2019 and 2020 (the later effect).

\paragraph{Cohen's kappa}
The raw agreement between two readers depends on the frequency of the answers: two readers who almost
always answer \file{no} often agree by chance. The kappa of \textcite{cohen1960coefficient} removes this
share:
\begin{equation}
  \kappa = \frac{p_o - p_e}{1 - p_e},
  \qquad
  p_e = \sum_c \pi_a(c)\, \pi_b(c),
\end{equation}
where $p_o$ is the share of identical answers and $\pi_a(c)$ the share of answers $c$ given by reader
$a$. The kappa is undefined when $p_e = 1$. The script \file{agree\_v2.py} does not compute it when one
of the two readers always gives the same answer. On ten papers, one or two disagreements drag the kappa
down towards 0.5 or 0.6 when a reader answers almost always the same way. The kappa is read
together with the share of agreement, the number of papers and its bootstrap interval
(Section~\ref{sec:sm-doubles}).

\paragraph{Cochran-Mantel-Haenszel test and Spearman correlation}
For the robustness analyses, a Cochran-Mantel-Haenszel test \parencite{mantel1959statistical} combines
the $2 \times 2$ tables of each conference (ICML, NeurIPS, ICLR), with a Haldane correction of 0.5 per
cell and a continuity correction (\file{statsmodels.stats.contingency\_tables.StratifiedTable}). The
Spearman correlation between the cost or risk scores and the observed differences receives a permutation
$p$-value: 10,000 permutations of the differences, NumPy generator with random seed 2026.

\section{The robustness analyses}\label{sec:sm-robustesse}

All these analyses are exploratory. None changes a verdict of deviation 20. They only say whether a
verdict is fragile.

\paragraph{Four cuts of the verdicts (deviation 21, a)}
For the four criteria that changed (Q5, Y9, X7, X2) and the two intervals that avoid zero without
passing Holm (Q6, Y12), \file{study/robustness.py} redoes the comparison in four ways.
\begin{itemize}
  \item \textbf{(a1) Provenance.} Classical RL in 2026 is split by provenance: the 14 papers of the
  first draw, with no title screen, against the 36 of the A3 queue, with the screen. The main
  comparison is redone with each provenance alone.
  \item \textbf{(a2) PDF source.} The 2026 papers are split by the source of the PDF read, arXiv or the
  conference site.
  \item \textbf{(a3) Conference.} A Cochran-Mantel-Haenszel test stratifies by conference, to rule out
  a composition effect.
  \item \textbf{(a4) Reader.} The comparison is redone removing each primary reader in turn. This cut
  does not see an effect shared by all the readers of one year (Section~\ref{sec:sm-annees}).
\end{itemize}
A verdict is called fragile if the difference changes sign in a cut, or if the recomputed Newcombe interval
contains zero while each side keeps at least 35 papers. Cut (a2) first read a damaged local copy of a
file. Redone on the repository file, it sets 45 papers read on arXiv against 5 read on a conference
site and settles nothing (deviation 21, c). For Q5, the re-reading of Section~\ref{sec:sm-versions}
measures the effect of the version read directly.

\paragraph{The cost and the risk of each practice (deviation 21, b)}
An earlier version of our conclusions claimed that Pineau had been heard where listening to her cost
little. That judgement about cost came from us, formed after seeing the differences. To put it to the test, three independent
raters each score the 17 practice criteria of the main comparison on two scales from 1 to 5: the cost
$C$ for an author who already has the main results, then the risk $R$ that the practice weakens or
reverses the conclusion. The raters are three AI agents of three different sizes, with no
access to the results or to the repository, who receive only the title and the rule of each criterion
(\file{study/blind\_cost/}). The test is the Spearman correlation between the median of the three scores
and the observed difference, with X3 reversed because a lower share is progress for it. A clear negative
correlation would support the earlier claim.

\paragraph{Without the thirteen papers read partly by another model (deviation 29)}
The thirteen papers finished by \file{claude-opus-5} are removed, then the main comparison and the
trajectory are redone with the same tests (environment variable \file{EXCLUDE}, outputs
\file{extension\_stats\_sans\_opus5.json} and \file{trajectory\_sans\_opus5.json}). A verdict changes if
a difference or a trend moves from one side of the threshold to the other after the Holm correction.

\paragraph{The title filter (A6, deviation 24)}
For each of the 400 papers drawn outside the filter, the abstract is read on the official page. A paper
is \enquote{hidden RL} if its abstract matches the same regular expression as the titles. The script
\file{study/title\_filter.py} gives the share of hidden RL with its Wilson interval, the estimated
number of papers missed (this share times the size of the population outside the filter) and their
fraction of the total. The first 20 hidden-RL papers in the order of the draw are read with both grids,
then compared descriptively with the 2026 sample on the four criteria that changed
(\file{study/title\_filter\_compare.py}). Four of them had no freely accessible PDF and were replaced by
the next ones in order.

\paragraph{The declarations of the papers (A5, deviation 22)}
NeurIPS 2025 papers attach the checklist filled in by their authors. The script \file{study/checklist.py}
retrieves the official proceedings version of each of the 30 papers by its title, then extracts the
answer and the justification of four questions: \enquote{Open access to data and code},
\enquote{Experimental setting/details}, \enquote{Experiment statistical significance},
\enquote{Experiments compute resources}. The comparison, fixed in advance, sets them against our
harmonised codes: Q5, Q3, X2 for an uncertainty and Q8 for a test, X7. The measure is the share of
papers that declare \enquote{Yes} while our code is \file{no}, with its Wilson interval. For each
disagreement on code, the script looks for a link to code in the proceedings version, because our codes
often come from the arXiv version.

\section{Two re-readings added after the results}\label{sec:sm-relectures}

A technical audit raised two points the earlier checks had left open: what X2
measures, then the version of the papers read for Q5. The two re-readings that answer them were written
into the protocol after the results and before any computation (deviations 32 and 31). They are
exploratory: no verdict of deviation 20 depends on them.

\paragraph{Criterion X8: an uncertainty computed over random seeds}\label{sec:sm-x8}
X2 counts any error bar (Section~\ref{sec:sm-grilles}). Yet a bar computed over the episodes of a
single run is the trap of 2018 that experiment B1 reconstructs (Section~\ref{sec:b1}). Deviation 32
defines criterion X8. It is \file{yes} when the uncertainty of the main results is computed across at
least two independent training runs and the paper says so (\enquote{mean and standard deviation over 5
seeds}). It is \file{no} when the uncertainty is computed over episodes, evaluation passes or test
items, when it mixes seeds and episodes in one spread, or when its source is not stated. For an
\file{llm} paper, only independent RL training runs count, never evaluation samples. Each reading also
fills in a \file{source} field: \file{seeds}, \file{pooled} (seeds and episodes mixed),
\file{episodes}, \file{items}, \file{unstated} or \file{other}.
\begin{itemize}
  \item \textbf{The sample.} The 74 papers of the main comparison whose final X2 is \file{yes}: 30
  from 2018 and 44 from classical RL in 2026. In the main measure, a paper with X2 \file{no} counts as
  \file{no} for X8, so the denominators stay the 48 and 50 papers with experiments.
  \item \textbf{The readers.} Eight new agents, Claude Code subagents on the \file{opus} model, which
  had read no paper of the survey, receive the instruction \file{study/x8\_prompt.txt}. It forbids them
  the other agents' readings and every result file. The sorted identifiers, shuffled by
  \file{random.Random(2032)}, are dealt out in turn among S1 to S6: each reads 12 or 13 papers from
  both years. S7 and S8 reread blind 20 papers drawn by \file{random.Random(2033)}, 10 per year, 5 and
  5 each (\file{study/x8\_assign.py}).
  \item \textbf{The check.} Each quotation is searched for in the PDF as for the other criteria
  (\file{study/x8\_check.py}). No file remained failing.
  \item \textbf{The analysis.} The script \file{study/x8\_stats.py} compares the two years with
  Newcombe and Fisher, alone and then with a Holm correction over the eleven v1 criteria plus X8.
  Beside it: the share among papers with X2 \file{yes}, the same measure with \file{pooled} counted
  \file{yes}, then a mechanical bound that searches the quotations of X2, Y2 and Y11 for a word naming a random
  seed or a run (deviation 30, e). The agreement of the double readings receives the intervals of
  Section~\ref{sec:sm-doubles} (generator with seed 2032).
\end{itemize}
X8 was read after the results, by agents other than the primary readers (result in
Section~\ref{sec:sr-principale}).

\paragraph{Q5 reread on the same kind of version}\label{sec:sm-versions}
Of the 50 classical RL papers from 2026, 45 were read on arXiv, where a link to code is often added after
acceptance. The papers of 2018 come mostly from the proceedings. Part of the rise of Q5 could therefore
come from the version read. Deviation 31 rereads Q5 on the same kind of version in both years,
mechanically, with the Q5 pattern of deviation 27 unchanged (\file{study/q5\_versions.py}).
\begin{itemize}
  \item \textbf{The versions.} For each of the 33 papers of 2018 where Q5 is \file{no}, the arXiv
  interface is queried with the title. The entry whose title matches at least 85\% gives the latest
  version: 27 are found, 6 do not exist on arXiv. For each of the 45 papers of 2026 read on arXiv, the
  proceedings version is retrieved: NeurIPS 2025 already on disk, ICLR 2026 on its proceedings site,
  ICML 2026 in volume 306 of PMLR by title. All 45 are found.
  \item \textbf{The decisions.} The text comes from \file{pdftotext}, with the references section
  removed. For a \file{yes} of 2026, the cited address is searched for in the proceedings version: if
  found, the decision stands. Otherwise, as for every \file{no}, each match of the pattern is kept with
  300 characters of context. The matches of both years are shuffled (\file{random.Random(31)}), the
  identifiers replaced by numbers, the years masked. A new arbiter judges them under the rules of
  \file{GRILLE.md}. An \enquote{uncertain} answer counts as \file{no} in the main reading and as
  \file{yes} in its bound.
  \item \textbf{The repository dates.} For each GitHub link of a \file{yes}, in either version, the
  creation date of the repository comes from the GitHub interface, one request at a time, 1.5 to
  4~seconds apart. It is compared with the last day of the conference.
  \item \textbf{The readings.} P0 keeps Q5 as published. P1 reads both years on the proceedings. P2
  reads them on the latest arXiv version. P3 keeps code public at the date of the conference: it is P2
  where a \file{yes} whose repositories were all created after the conference becomes \file{no}. Each
  reading redoes the main comparison with the same test: Newcombe, Fisher, Holm over the eleven v1
  criteria with the ten other $p$-values unchanged (\file{study/q5\_versions\_stats.py}).
\end{itemize}
The downloads targeted one host at a time, 3 to 6~seconds apart for arXiv. No request was refused.
Section~\ref{sec:sr-robustesse} gives the seven readings.

\section{Deviations from the protocol}\label{sec:sm-ecarts}

The survey protocol counts 35 deviations, numbered 1 to 36: number 34 awaits a past human re-reading and will be written only with its facts. Table~\ref{tab:ecarts} gives all of them with their date and
their position relative to the results, next to those of the two other protocols
(Chapter~\ref{chap:09-ecarts}). They fall into five groups.

\begin{itemize}
  \item \textbf{Rules clarified before any percentage}: deviations 1 to 6 (grid v1), 14, 16 to 19
  (grid v2, pilot, X5 family, harmonisation). They remove ambiguities without targeting a number.
  \item \textbf{Instruments added before their own computation}: deviations 11 and 12 (grid v2,
  enlargement), 13 and 23 (human rereading planned, then abandoned before it started), 22, 24, 25 and 26 (A5, A6, A4 and its
  second reading).
  \item \textbf{Decisions taken after seeing results}: deviations 7 (initial calibration), 15 (2018
  harmonisation caught up), 20 (same-grid verdict rule), 21, 27, 28 and 29 (analyses answering a
  critical review), 30 to 33 (response to the technical audit), 35 and 36 (human rereading of both
  years, eight fresh readers on both years). The analyses were all
  written before their own computation. The protocol declares these decisions as such.
  \item \textbf{Incidents with no effect on the data}: deviations 9 (contact address removed from the
  download header) and 10 (working scripts overwritten in a shared temporary folder, codes intact).
  \item \textbf{Reading facts}: deviation 8 (appendices of 2018 and 2026).
\end{itemize}
